\documentclass[11pt]{letter}

\usepackage[margin=0.75in]{geometry}
\usepackage{setspace}
\usepackage{hyperref}
\usepackage{xcolor}
\usepackage{enumitem}
\usepackage{parskip}

\hypersetup{
    colorlinks=true,
    urlcolor=black,
    linkcolor=black
}

\setlength{\parindent}{0pt}
\setstretch{1.05}
\date{}
\begin{document}

\begin{letter}{}

\opening{Dear Hiring Committee,}

I saw the job post on the \textbf{Notion} careers site and thought I'd be a good fit for it.

I am a Master's student in Computer Science at the University of Cincinnati, graduating December 2026, looking for a role as a \textbf{Software Engineer, New Grad} at \textbf{Notion}. With hands on experience building web applications, backend services, and reliable deployment systems, I am excited about the opportunity to contribute to \textbf{Notion} and support the team building the collaborative AI workspace millions of people rely on.

\textbf{Here are reasons that make me a great fit for the role:}

\begin{itemize}[leftmargin=*, itemsep=0.45em]
    \item \textbf{Full-stack web application experience:} I built PolicyPilot, a full-stack TypeScript application with a React frontend and Express/Node.js backend, including a 17-tool MCP server with Zod schema validation at every data boundary. I also built a React and MySQL workflow application for the Government of Maharashtra that digitized departmental record management and improved operational efficiency by 40\%.

    \item \textbf{Familiarity with backend systems and databases:} On the Spotify Personal Analytics Platform, I designed a PostgreSQL analytics schema with materialized views and 16 SQL RPC functions, and built a FastAPI backend serving low-latency queries over 71,000+ streaming events, giving me hands on experience connecting the data model to the user interface.

    \item \textbf{Pragmatic use of AI-assisted development:} In my current graduate research, I designed a memory-augmented LLM pipeline for Android bug reproduction using Gemini and Vertex AI, achieving a 25.7x latency reduction over a published baseline. I regularly use AI-assisted development tools like Claude Code in my own workflow to move faster while staying focused on the problem I am solving.

    \item \textbf{Production engineering and reliability mindset:} At Siemens Industry Software, I automated infrastructure workflows with Python and Boto3, cutting deployment time by 40\%, and executed 44 production environment upgrades in 4 months while diagnosing live issues across Linux, EC2, and Kubernetes services. As Quality Coach and SonarQube Specialist, I increased automated code coverage from 35\% to 78\% and mentored 5 engineers.
\end{itemize}

I am available to start full time after my graduation in December 2026. As a Master's student at the University of Cincinnati and someone who values professionalism and growth, I am enthusiastic about contributing to Notion's mission of making work faster, clearer, and less fragmented.

Please find my CV attached below. I look forward to hearing from you. I would be excited to help build tools that let humans and AI agents think and work together at Notion.

{\setlength{\parskip}{0pt}
Regards,\\
Tanmay Bhuskute,\\
tanmay.v.bhuskute@gmail.com $\mid$ +1 (732) 858-3027,\\
San Francisco, CA.
}

\end{letter}

\end{document}
